% 4-2-hourly.tex
%  Budget: 6pp.  The canonical results table is the spine of 4-2 and 4-3.
%  Source map: the published protocol and reference numbers this table is
%   built to be comparable with -> lago_forecasting_2021 .
%  All values transcribed at 4 dp from reports/tables/results_canonical.csv
%   (v1.0-results, frozen). Hourly rows:
%    naive     7.7501 13.2573 28.5954 0.8491
%    SARIMAX   4.3507  7.1166 18.0349 0.4767
%    LEAR      3.8988  6.4750 16.6572 0.4272
%    LightGBM  3.9678  7.5024 15.7218 0.4347
%    LSTM      3.8734  7.0198 15.5994 0.4244
%    static    3.5742  6.6097 14.6710 0.3916
%    regime    3.5569  6.5575 14.6439 0.3897
%  Table numerals are Latin inside an RTL tabular -- verified on render;
%  Persian decimal marks are not in B Nazanin (see fix_numerals.py).

\section{نتایج پیش‌بینی ساعتی}
\label{sec:4-2}

این بخش دقتِ پیش‌بینیِ بردار ۲۴مؤلفه‌ای قیمت‌های ساعتی روز بعد را گزارش می‌کند: هدف
اصلی پژوهش، و پایه‌ای که نتایج بخش‌های بعدی بر آن بنا می‌شوند. تمام اعداد بر ۷۲۸ مبدأ
پیش‌بینی و \lr{17{,}472} ساعتِ دوره‌ی آزمون محاسبه شده‌اند.

\begin{table}[htbp]
\centering
\caption{دقت پیش‌بینی ساعتی بر دوره‌ی آزمون \lr{EPEX-DE}، از ۴ ژانویه‌ی ۲۰۱۶ تا ۳۱
دسامبر ۲۰۱۷ (۷۲۸ مبدأ). واحد \lr{MAE} و \lr{RMSE} یورو بر مگاوات‌ساعت، \lr{sMAPE}
درصد، و \lr{rMAE} بی‌بعد. مقدار کمتر بهتر است. مدل ساده مرجع است و نه رقیب.}
\label{tab:hourly-results}
\small
\setlength{\tabcolsep}{6pt}
\begin{tabular}{lrrrr}
\toprule
مدل & \lr{MAE} & \lr{RMSE} & \lr{sMAPE} & \lr{rMAE} \\
\midrule
ساده & \lr{7.7501} & \lr{13.2573} & \lr{28.5954} & \lr{0.8491} \\
\lr{SARIMAX} & \lr{4.3507} & \lr{7.1166} & \lr{18.0349} & \lr{0.4767} \\
\lr{LEAR-LASSO} & \lr{3.8988} & \lr{\textbf{6.4750}} & \lr{16.6572} & \lr{0.4272} \\
\lr{LightGBM} & \lr{3.9678} & \lr{7.5024} & \lr{15.7218} & \lr{0.4347} \\
\lr{LSTM} & \lr{3.8734} & \lr{7.0198} & \lr{15.5994} & \lr{0.4244} \\
\midrule
ترکیب ایستا & \lr{3.5742} & \lr{6.6097} & \lr{14.6710} & \lr{0.3916} \\
ترکیب رژیم‌محور & \lr{\textbf{3.5569}} & \lr{6.5575} & \lr{\textbf{14.6439}} & \lr{\textbf{0.3897}} \\
\bottomrule
\end{tabular}
\end{table}

\subsection*{ترتیب مدل‌ها}

بر معیار \lr{MAE}، ترتیب چنین است: ترکیب رژیم‌محور با \lr{3.5569}، ترکیب ایستا با
\lr{3.5742}، \lr{LSTM} با \lr{3.8734}، \lr{LEAR-LASSO} با \lr{3.8988}،
\lr{LightGBM} با \lr{3.9678}، \lr{SARIMAX} با \lr{4.3507} و مدل ساده با
\lr{7.7501}.

سه مشاهده از این ترتیب برمی‌خیزد که هر سه در \cref{sec:4-5-1} با آزمون معناداری سنجیده
می‌شوند؛ در اینجا تنها به‌عنوان مشاهده‌ی عددی ثبت می‌گردند.

\textbf{یکم: فاصله‌ی میان سه مدل میانی بسیار کوچک است.} \lr{LSTM}، \lr{LEAR-LASSO} و
\lr{LightGBM} در بازه‌ای به پهنای حدود هفت صدم یورو بر مگاوات‌ساعت جای می‌گیرند
(\lr{3.8734} تا \lr{3.9678}). در برابر انحراف معیارِ خودِ سری قیمت، که \lr{15.94}
است، این فاصله ناچیز است. \cref{sec:4-5-1} نشان می‌دهد که هیچ‌یک از تفاوت‌های درون این سه
مدل، پس از تصحیح برای چندگانگی، معنادار نیست.

\textbf{دوم: مدل خطی از مدل درختی جلوتر است.} \lr{LEAR-LASSO} با \lr{3.8988} از
\lr{LightGBM} با \lr{3.9678} بهتر عمل می‌کند. این نتیجه‌ای است که انتظار متعارف را
برآورده نمی‌کند و باید بی‌تعدیل گزارش شود: در این بازار و این دوره، انعطاف‌پذیریِ
مدل درختی به دقت بیشتر ترجمه نشد.

\textbf{سوم: هر دو ترکیب از تمام اعضایشان بهترند.} کوچک‌ترین \lr{MAE} در میان
تک‌مدل‌ها \lr{3.8734} است، حال آنکه ترکیب ایستا به \lr{3.5742} و ترکیب رژیم‌محور به
\lr{3.5569} می‌رسد. این بهبود از هر بهبودی که میان خودِ تک‌مدل‌ها دیده می‌شود بزرگ‌تر
است، و برخلاف آن‌ها از آزمون معناداری نیز عبور می‌کند.

\subsection*{معیارها همیشه هم‌داستان نیستند}

\cref{tab:hourly-results} یک ناسازگاریِ آموزنده دارد که پوشاندنش آسان و
گزارش‌کردنش درست‌تر است: بهترین مدل بر معیار \lr{RMSE}، ترکیب رژیم‌محور نیست، بلکه
\lr{LEAR-LASSO} با \lr{6.4750} است.

توضیح در تفاوت این دو معیار نهفته است. \lr{RMSE} خطاها را به توان دو می‌رساند و در
نتیجه به خطاهای بزرگ حساسیت به‌مراتب بیشتری دارد تا \lr{MAE}. یک مدل می‌تواند
میانگین خطای مطلقِ بهتری داشته باشد و هم‌زمان چند خطای بسیار بزرگ‌تر مرتکب شود.
مقایسه‌ی \lr{LEAR-LASSO} و \lr{LightGBM} همین را نشان می‌دهد: مدل درختی در \lr{sMAPE}
بهتر است (\lr{15.7218} در برابر \lr{16.6572}) اما در \lr{RMSE} به‌روشنی بدتر
(\lr{7.5024} در برابر \lr{6.4750}) \lr{—} یعنی خطاهای معمولیِ کوچک‌تر اما خطاهای
حدیِ بزرگ‌تر، که با ماهیت درختی و اشباع‌شدن آن در دنباله‌ها سازگار است.

این ناسازگاری بر هیچ نتیجه‌ای در این پایان‌نامه اثر نمی‌گذارد، چون \lr{—} همان‌گونه
که \cref{sec:3-5} تصریح کرد \lr{—} تمام آزمون‌های معناداری بر زیان قدرمطلق محاسبه می‌شوند و
هر ادعای سرشناسی بر \lr{MAE} یا \lr{rMAE} بنا شده است. اما ثبت‌کردنش لازم است: انتخاب
معیار، ترتیب مدل‌ها را تغییر می‌دهد، و پژوهشی که تنها معیارِ مساعد به حال خود را
گزارش کند، همین واقعیت را پنهان کرده است.

\subsection*{خواندن معیار \lr{rMAE}}

ستون \lr{rMAE} همان اطلاعات \lr{MAE} را در مقیاسی تفسیرپذیر می‌دهد: نسبت خطای مدل به
خطای پیش‌بینیِ ساده‌ی وقفه‌ی هفتگی. مقدار \lr{0.3897} برای ترکیب رژیم‌محور یعنی این مدل
حدود شصت‌ویک درصد از خطای آن مرجع را حذف می‌کند.

مقدار \lr{0.8491} برای مدل ساده نیز، همان‌گونه که \cref{sec:3-5} توضیح داد، تناقض نیست: مدل
ساده‌ی گزارش‌شده قاعده‌ی روز مشابه است، حال آنکه مخرجِ \lr{rMAE} پیش‌بینیِ وقفه‌ی هفتگیِ
محض است، و قاعده‌ی روز مشابه حدود پانزده درصد از آن بهتر عمل می‌کند.

\subsection*{آنچه این بخش ادعا نمی‌کند}

اعداد این جدول به یک بازار، یک پنجره‌ی دوساله، یک معماری به‌ازای هر خانواده‌ی مدل و یک
بودجه‌ی تنظیمِ پنجاه‌آزمایشی محدودند. این بخش نمی‌گوید که یادگیری عمیق بر روش‌های آماری
برتری دارد یا ندارد؛ می‌گوید که در این شرایط مشخص، اختلاف آن‌ها کوچک بود. تعمیم این
مشاهده به یک حکم کلی درباره‌ی خانواده‌های مدل، فراتر از شواهدی است که در دست است.

مقایسه با اعداد منتشرشده‌ی مرجع در \cref{sec:4-5-2} می‌آید، و پرسشِ اینکه کدام‌یک از این
تفاوت‌ها معنادارند در \cref{sec:4-5-1} پاسخ داده می‌شود.
